% ==================================================================
% id:       sets-004
% topic:    תורת הקבוצות ומאורעות
% subtopic: מכפלה קרטזית וקבוצת החזקה
% solution: full
% ==================================================================
יהיו $A=\set{1,2,3}$ ו-$B=\set{a,b}$.
\begin{enumerate}
\item כתבו את $A\times B$ ואת $B\times A$. האם הן שוות?
\begin{solbox}
לפי ההגדרה,
\[
A\times B=\set{(1,a),(1,b),(2,a),(2,b),(3,a),(3,b)},\qquad
B\times A=\set{(a,1),(a,2),(a,3),(b,1),(b,2),(b,3)}.
\]
בשתי הקבוצות יש 6 איברים, אבל הן \textbf{אינן שוות}: למשל $(1,a)\in A\times B$, אבל $(1,a)\notin B\times A$, כי בזוג סדור יש חשיבות לסדר, ו-$(1,a)\neq(a,1)$. למעשה אין להן אף איבר משותף.
\end{solbox}
\item כמה איברים יש ב-$\Pow(A\times B)$?
\begin{solbox}
לפי עקרון הכפל, $|A\times B|=|A|\cdot|B|=3\cdot2=6$. לקבוצה עם 6 איברים יש $2^6$ תתי-קבוצות, ולכן $|\Pow(A\times B)|=2^6=\ans{64}$.
\end{solbox}
\end{enumerate}
